% conclusion.tex
%
% Purpose:
%   Close the first-draft report and identify near-term refinements.
%
% Main algorithm:
%   No algorithm is introduced.  The section summarizes the formulation and
%   lists open documentation and validation tasks.
%
% Based on:
%   The preceding sections of this report.
%
% Main changes:
%   Provides a modest conclusion suitable for a draft technical note.
%
% Numerical goal:
%   Make clear what should be checked next before treating the report as a
%   polished paper manuscript.

\section{Conclusion}

This report has organized the current Transmission Eigenvalue / Bloch-mode /
line-defect formulation into a single technical narrative.  The main numerical
structure is a center-cell boundary integral equation coupled to compact port
matching conditions.  The half-lead exterior condition is imposed by requiring
the port Cauchy data to lie in Bloch-selected incoming or outgoing trace
subspaces, rather than by explicitly forming a full periodic half-guide
Dirichlet-to-Neumann operator.

The formulation is deliberately practical.  It identifies which pieces are
continuous modeling assumptions, which pieces are finite-dimensional
approximations, and which choices are numerical heuristics that need convergence
checks.  The most important next steps are to insert validated numerical
examples, record parameter choices for stable scans, and compare successive
refinements of \(N_{\mathrm{tot}}\), \(M\), proxy data, and mode-selection
thresholds.

Several details remain draft-level.  The exact treatment of unit-circle Bloch
modes should be documented more fully once the energy-flux selection is settled.
The bibliography metadata should be checked against the final source PDFs or
publisher records.  Finally, the report should eventually include numerical
figures showing singular-value scans, field plots, and convergence tables for
representative line-defect configurations.
